\begin{figure*}[!t]
  \centering
  \includegraphics[width=\textwidth]{\paperroot figures/fusion-architecture-en.pdf}
  \caption{Quality-gated fusion and its safety boundary. Solid arrows show the
  current workflow: 19 features pass through modality-specific quality gates
  and combine with quality and missingness fields into a 29-dimensional
  research representation. In parallel, the five component statuses produced
  by their upstream modality-specific rules enter the lower safety path, which
  determines user-facing urgency independently of the research vector. The
  dashed branch marks a probability-model extension based on clinically labeled
  subjects, calibration, and external evaluation.}
  \label{fig:fusion_architecture}
\end{figure*}

\section{Quality-Gated Multimodal Fusion}
\label{sec:fusion}

The fusion layer converts completed B/E/F/A/S measurements into a fixed research
representation. It contains 19 continuous features, five modality-quality
values, and five modality-missingness indicators, for a total of 29 dimensions;
these counts come from the implementation's fixed feature list, the five
BE-FAST modalities, and their per-modality missingness flags. The continuous
subsets contain 2, 6, 3, 4, and 4 features for B, E, F, A, and S, respectively,
as summarized in Fig.~\ref{fig:fusion_architecture}. Let
$\mathbf{x}_m$ be the continuous feature subset and $q_m\in[0,1]$ the quality
value for modality $m$. The module-reported valid-frame fraction or audio-quality
support initializes $q_m$; where action completion applies, it is the minimum of
that support and completion terms clipped to $[0,1]$. Thus, face quality also
depends on $a_{\max}/0.22$, and arm quality also depends on completing the full
action sequence. Quality-gated features are

\begin{equation}
\tilde{\mathbf{x}}_m=q_m\mathbf{x}_m.
\label{eq:quality_gating}
\end{equation}

The research vector is

\begin{equation}
\mathbf{z}=[\tilde{\mathbf{x}}_B,\ldots,\tilde{\mathbf{x}}_S,
q_B,\ldots,q_S,r_B,\ldots,r_S],
\label{eq:fusion_vector}
\end{equation}

where $r_m=1$ denotes a missing modality. Retaining $q_m$ and $r_m$ distinguishes
a zero-valued measurement from an unavailable observation. The fusion layer
also produces domain severities, a top-two summary, and arm--face laterality
agreement for domain comparison and multimodal research analysis.

For each modality, a domain severity is computed symmetrically from its
quality value and its available reference-normalized feature deviations:

\begin{equation}
d_m=q_m\min\!\left\{c,
\max_{j\in\mathcal{J}_m}g_{m,j}(x_{m,j})\right\},
\label{eq:domain_fusion_severity}
\end{equation}

where $\mathcal{J}_m$ indexes the modality's implemented severity features,
$g_{m,j}$ maps a measurement to a nonnegative deviation relative to its
prespecified engineering reference, and $c$ clips extreme values. These
values form the modality-severity vector
$\mathbf{d}=[d_B,d_E,d_F,d_A,d_S]$, from which the maximum domain severity and
the mean of the two largest domain severities are calculated.

As shown in Fig.~\ref{fig:fusion_architecture}, each B/E/F/A/S subsystem also
produces a component status upstream using the modality-specific rules defined
in Section~\ref{sec:methods}. An independent safety path combines those statuses
with onset information: a positive finding with new or sudden onset is
escalated; unclear onset produces a warning; valid negative results across all
requested tasks produce a clear outcome; and insufficient or skipped tasks mark
the screen incomplete. Future work will extend the research representation to a
clinically interpretable probability model by prespecifying the prediction
target and reference standard, using participant-disjoint development and
evaluation data, and assessing discrimination, calibration, and clinical
utility through representative external validation \cite{collins2024tripodai}.
